一个迭代算法跑了很久，输出的数字越来越稳定。它收敛了吗？这句话里藏着一个被默认掉的前提：它想去的那个点，得在你研究的空间里真实存在。有理数系统给不出这个保证——一列有理数可以无限逼近 \(\sqrt2\)，却没有一个有理数等于它。极限、连续、积分全都建立在``逼近不会掉进空洞''之上，所以在讲这些之前，先得把空洞补上。

本章用两种语言说同一件事。上确界公理从集合的边界入手：非空且有上界的集合必有最小上界，那个最小上界就落在原本空缺的位置上。Cauchy 判据从序列的内部入手：只要尾部各项彼此挤在一起，极限就存在——注意这句话在 \(\Q\) 里是假的，它成立与否恰好等价于完备性。

一个具体的定位点值得记住。在 \(\Q\) 中取 \(A=\set{q>0:q^2<2}\)，找它的最小上界会落空；取逼近序列 \(1,\,1.4,\,1.41,\,1.414,\ldots\)，各项彼此无限接近却无处可归。两条路走到同一个洞口，补洞的方式也因此有两套等价的说法。后面反复出现的单调有界收敛、闭区间套、Bolzano--Weierstrass，用的都是这一条资产的不同支取方式。

\subsection{怎么读这一章}

两篇都在回答同一个问题：逼近为什么不会掉出数系之外。建议先弄清它们各自解决的困境，再去记等价表述——记住``上确界公理等价于 Cauchy 完备''这句话很容易，看出每次写下 \(\sup\) 或每次用相邻差判停时都在动用它，才是这一章的用处。

\subsection{本章篇目}

\begin{itemize}
\tightlist
\item \hyperref[sec:12-Real-Analysis-Support/01-Real-Numbers-and-Completeness/01]{``上确界填补有理数的空洞''}：稠密不等于没有洞，最小上界公理如何一次性补齐，以及为什么单调有界收敛的第一行字依赖它；
\item \hyperref[sec:12-Real-Analysis-Support/01-Real-Numbers-and-Completeness/02]{``Cauchy 判据无需预知极限''}：把参照物从极限换成序列自身的尾部，以及``相邻差趋于零''为什么不够。
\end{itemize}

读完这两篇，你会得到一个可以随身携带的检查动作：每当一段论证写出``令极限为 \(L\)''，先问一句 \(L\) 凭什么存在、存在于哪个空间。这个问题在 \(\R\) 上有现成答案，换到函数空间就需要重新回答，那是下一章的事。
